%------------------------------------------------------------------------------%
\documentclass[a4paper,12pt,fleqn]{article}

\usepackage{calculo_varias_variaveis-1}

\ShowAnswers

%------------------------------------------------------------------------------%
\begin{document}

\makeHeader{CFVVI}{Prova 2 -- Turma 4}{04/06/2025}

% 01
%------------------------------------------------------------------------------%
\exercise{25}
Considere que a equação
\(
  e^{xyz} + x^2 - z = 0
\)
defina $z$ como função de $x$ e $y$, encontre \(\D{z}{x}\)
\clearpagequestiononly

\begin{answer}
  \begin{align*}
    e^{xyz(x, y)} + x^2 - z(x, y)                                    & = 0 \\
    \D{}{x}\left( e^{xyz(x, y)} + x^2 - z(x, y) \right)              & = 0 \\
    \D{}{x}e^{xyz(x, y)} + \D{x^2}{x} - \D{z}{x}                     & = 0 \\
    e^{xyz(x, y)}\D{}{x}\left(xyz(x, y)\right) + 2x - \D{z}{x}       & = 0 \\
    e^{xyz(x, y)}\left(yz(x, y) + xy \D{z}{x}\right) + 2x - \D{z}{x} & = 0 \\
    \left( xy e^{xyz(x, y)} - 1 \right) \D{z}{x}
    & = - yz(x, y) e^{xyz(x, y)} - 2x  \\
    \D{z}{x}
    & = \frac{ - yz(x, y) e^{xyz(x, y)} - 2x }{ xy e^{xyz(x, y)} - 1 } \\
    & = \frac{ - yz e^{xyz} - 2x }{ xy e^{xyz} - 1 }
  \end{align*}
\end{answer}

% 02
%------------------------------------------------------------------------------%
\exercise{25}
Seja \( f(x, y) = \ln(x^2 + y^2) \),
use a aproximação linear de $f$ no ponto $(1, 2)$
para estimar $f(1,01; 1,98)$
\clearpagequestiononly

\begin{answer}
  Calculando as derivadas parciais
  \begin{align*}
    \D{f}{x} & = \frac{2x}{x^2 + y^2} \\
    \D{f}{y} & = \frac{2y}{x^2 + y^2}
  \end{align*}
  Avaliando as derivadas no ponto $(1, 2)$
  \begin{align*}
    f(1, 2)        & = \ln(1 + 4) = \ln(5) \\
    \D{f}{x}(1, 2) & = \frac{2}{5}         \\
    \D{f}{y}(1, 2) & = \frac{4}{5}
  \end{align*}
  A aproximação linear é
  \begin{align*}
    L(x, y) & = f(x_0, y_0) + \D{f}{x}(x_0, y_0)(x-x_0) + \D{f}{y}(y-y_0) \\
            & = \ln(5) + \frac{2}{5}(x - 1) + \frac{4}{5}(y - 2)
  \end{align*}
  Avaliando no ponto $f(1,01; 1,98)$
  \begin{align*}
    L(1,01; 1,98) & = \ln(5) + \frac{2}{5}(0,01) + \frac{4}{5}(-0,02) \\
                  & = \ln(5) + \frac{2-8}{500} \\
                  & = \ln(5) - \frac{6}{500} \\
                  & = \ln(5) - \frac{3}{250} \\
                  & \approx 1.59744
  \end{align*}
  Comparando com o valor exato
  \[
    f(1,01; 1,98) = 1.59747
  \]
\end{answer}

% 03
%------------------------------------------------------------------------------%
\exercise{25}
Seja \( f(x, y) = \ln(x^2 + y^2 + 1) \)
\begin{enumerate}[label=\alph*)]
  \item Calcule o gradiente de $f$
  \item Calcule a derivada de $f$ no ponto $(1, 2)$ na direção do vetor $\vetor{-2}{1}$
\end{enumerate}
\clearpagequestiononly

\begin{answer}
  \begin{enumerate}[label=\alph*)]
    \item
      Vetor gradiente
      \[
        \nabla f(x, y)
        = \Vetor{\D{f}{x}}{\D{f}{y}}
        = \Vetor{\frac{2x}{x^2 + y^2 + 1}}{\frac{2y}{x^2 + y^2 + 1}}
      \]
    \item
      Gradiente no ponto $(1, 2)$
      \[
        \nabla f(x, y) (1, 2)
        = \Vetor{\frac{2}{6}}{\frac{4}{6}}
        = \Vetor{\frac{1}{3}}{\frac{2}{3}}
      \]
      Calculando um vetor unitário na direção de $v$
      \[
        u
        = \frac{v}{\norm{v}}
        = \frac{1}{\sqrt{(-2)^2 + 1^2}} \vetor{-2}{1}
        = \frac{1}{\sqrt{5}} \vetor{-2}{1}
      \]
      Derivada direcional
      \begin{align*}
        D_u f(1, 2)
        & = \nabla f(1, 2) \cdot u                       \\
        & = \Vetor{\frac{1}{3}}{\frac{2}{3}} \cdot \frac{1}{\sqrt{5}} \vetor{-2}{1} \\
        & = \frac{1}{\sqrt{5}} \left( \frac{-2}{3} + \frac{2}{3} \right) \\
        &  = 0
      \end{align*}
  \end{enumerate}
\end{answer}

% 04
%------------------------------------------------------------------------------%
\exercise{25}
Seja \( f(x, y, z) = e^{x+y} \cos(z) \),
calcule \( \DM{f}{x}{z} \) no ponto \( (0, 0, \pi) \).
Efetue as derivadas na ordem especificada pela notação.
\clearpagequestiononly

\begin{answer}
\begin{align*}
  \D{f}{z}
  & = \D{}{z}\left(e^{x+y} \cos(z)\right) \\
  & = e^{x+y} \D{}{z}\cos(z)\\
  & = e^{x+y} (-\sin(z)) \\
  & = -e^{x+y} \sin(z)
\end{align*}
\begin{align*}
  \DM{f}{x}{z}
  & = \D{}{x}\left(-e^{x+y} \sin(z)\right) \\
  & = -\sin(z)\D{}{x}e^{x+y} \\
  & = -\sin(z)e^{x+y}\D{}{x}\left(x+y\right) \\
  & = -\sin(z) e^{x+y}
\end{align*}
\begin{align*}
  \DM{f}{x}{z}(0, 0, \pi)
  & = -\sin(z) e^{x+y} \evalat{(0, 0, \pi)}{} \\
  & = -\sin(\pi) e^{0+0} \\
  & = -\sin(\pi) e^0 \\
  & = 0
\end{align*}
\end{answer}

%------------------------------------------------------------------------------%
\end{document}
%------------------------------------------------------------------------------%
